\documentclass[project_master.tex]{subfiles}
\begin{document}
\section{Project Objective}
The objective is to construct a compact, physics-based framework that connects radiation transport in silicon to semiconductor degradation. Geant4 provides the front-end description of particle interactions and energy transfer, while MATLAB provides the back-end defect and device-degradation model.

\section{Main Modules}
\subsection{Module 1: Geant4 Radiation Transport}
A simple silicon slab geometry is irradiated with protons or neutrons. The goal is not to build a full mission-grade transport model, but to obtain interpretable observables such as deposited energy, secondary production, recoil events, and depth-dependent damage proxies.

\subsection{Module 2: Defect Build-Up Model}
A first-order defect model maps fluence or a displacement-damage proxy to an effective defect concentration:
\begin{equation}
\Nd = k_D \PhiF,
\end{equation}
or, if depth information is retained,
\begin{equation}
\Nd(z) = k_D\, \dd(z).
\end{equation}

\subsection{Module 3: Lifetime and Transport Degradation}
Defects are treated as recombination-active centers that reduce the effective carrier lifetime:
\begin{equation}
\frac{1}{\teff} = \frac{1}{\tzero} + K_\tau \Nd.
\end{equation}
The corresponding diffusion length is then
\begin{equation}
L = \sqrt{D\,\teff}.
\end{equation}

\subsection{Module 4: Device-Level Degradation}
For a silicon diode or detector, leakage current growth can be modeled with
\begin{equation}
\Delta I = \alpha \PhiF V,
\end{equation}
where $V$ is the active or depleted volume.

\section{Deliverables}
\begin{itemize}[leftmargin=1.5em]
    \item Geant4 figures showing deposited-energy and depth-dependent damage observables.
    \item MATLAB figures showing defect concentration, lifetime, diffusion length, and leakage current versus fluence.
    \item A concise interview presentation that explains the physical chain from microscopic damage to measurable device degradation.
\end{itemize}

\section{Milestones}
\begin{enumerate}[leftmargin=1.5em]
    \item Finalize the conceptual model and notes.
    \item Build the MATLAB back-end first to establish visible project progress.
    \item Add the Geant4 front-end and connect its output to the MATLAB pipeline.
    \item Convert the finished work into slides and speaking notes.
\end{enumerate}

\section{Why This Project Works Well}
This project is attractive because it is technically broad without becoming unmanageably large. It shows familiarity with radiation-matter interaction, defect physics in crystalline silicon, and engineering interpretation at the device level. It also has a clear computational split: Geant4 handles transport, while MATLAB handles solid-state interpretation.
\end{document}
